\documentclass[11pt]{article}
\usepackage[margin=1in]{geometry}
\usepackage[T1]{fontenc}
\usepackage{amsmath,amssymb}
\usepackage{booktabs}
\usepackage[hidelinks]{hyperref}
\usepackage[expansion=false]{microtype}

\title{Trace Transport: A Self-Play Testbed for\\History-Based Learned Network Protocols}
\author{Project working draft}
\date{\today}

\begin{document}
\maketitle

\begin{abstract}
We present a testbed for exact text and binary-object delivery over datagram
networks with unknown and changing capacity, latency, loss, and competing hosts.
Both endpoints are controlled, allowing custom feedback and metadata without TCP
wire compatibility. A deterministic reliability kernel enforces byte correctness;
local history transformers control sender pacing, flight window, scheduling, and
receiver feedback. CPU event workers support joint MAPPO with batched GPU
learning and a training-only centralized critic. Deterministic feedback-budget
changes improve earlier diagnostic completion time, but learning has not shown
general superiority. In the mixed-workload joint study, six five-million-decision
runs finish and validation selects the initial policy. On 40 held-out sessions
its deadline-miss rate is 40.00\%, versus 40.71\% for the selected deterministic
reference, failing the predeclared 20\% relative-improvement target. The selected
policy is untrained, so this difference is not a learning gain. Thirty-six bounded
real-device transfers recover exact synthetic text, image, and map bytes; the
learned-sender checks retain deterministic remote feedback and do not establish
performance superiority. We document negative results, runtime limitations, and a
separate planned objective for maximum verified goodput with backlogged traffic.
\end{abstract}

\section{Motivation and scope}
An endpoint cannot directly observe network capacity, competing senders, router
queues, or future failures. Its decisions must depend on local transmission
records and received messages. A history-based policy can attempt to infer
latent conditions and adapt its communication behavior. Our concrete application
is exact text transfer: successful output must match the original byte sequence,
including whitespace and Unicode encoding. The policy never generates substitute
payload content.

Both endpoints run the experimental protocol. We can therefore optimize sender
and receiver behavior together, including the representation and timing of
feedback. The network remains an external datagram service. IP routing and
compatibility with unmodified TCP peers are outside the initial scope. Permanent
partitions produce explicit incomplete transfers; resumable sessions initially
assume retained endpoint state.

\section{Related work and positioning}
Remy synthesizes endpoint congestion-control rules from network assumptions and
an optimization objective~\cite{winstein2013remy}. Aurora studies deep reinforcement
learning for congestion control with a Gym-based experimental
interface~\cite{jay2019aurora}. MVFST-RL integrates asynchronous reinforcement
learning with QUIC and addresses action and inference delays in an operational
transport~\cite{sivakumar2019mvfst}.

Transformers over network histories also have direct precedent. The Jazz preprint
uses historical observations for transformer-based multipath congestion
control~\cite{wang2025jazz}. The TCCO preprint investigates transformer-based
control across wireless subflows~\cite{zhang2026tcco}. These works motivate a
history-based controller but mean that this architectural choice alone is not a
novel contribution.

Trace Transport's intended experimental focus is a compact combination of exact
application delivery, jointly controlled peers, budgeted custom feedback, and
scalable packet simulation. Whether this constitutes a useful contribution
requires implementation, comparison, and further literature review. TCP's
reliability specification and QUIC recovery provide reference mechanisms, not
mandatory wire formats~\cite{rfc9293,rfc9002}.

\section{Problem formulation}
Let $s_t$ denote the full network and endpoint state at control step $t$. This
includes packet queues, scheduled arrivals, hidden link conditions, and endpoint
buffers. Endpoint $i$ receives local observation $o_t^i$ and maintains history
\begin{equation}
 h_t^i = \bigl((o_\tau^i,a_{\tau-1}^i,\Delta t_\tau)\bigr)_{\tau=t-H+1}^{t}.
\end{equation}
Padding and validity masks handle startup and missing measurements. A shared
role-conditioned policy samples $a_t^i\sim\pi_\theta(\cdot\mid h_t^i)$.
Endpoints interact through network packets; sharing policy weights does not
give them shared runtime observations. Actions precede the next environment
transition. The task is partially observable and becomes a cooperative multi-agent
problem when both peers and multiple connections are trained jointly.

For source bytes $x$ and reconstructed output $\hat{x}$, verified success requires
\begin{equation}
 \operatorname{complete}(x,\hat{x}) \iff
 |x|=|\hat{x}|\ \land\ \forall j,\ x_j=\hat{x}_j.
\end{equation}
The verifier operates outside the policy. A deadline bounds finite-transfer
episodes; lack of connectivity does not imply that delivery is possible.

\section{Protocol and controller}
\subsection{Reliability kernel and metadata}
The initial kernel uses connection and message identifiers, byte offsets, final
lengths, unique transmission identifiers, explicit received ranges, reassembly,
and duplicate suppression. Integrity and completion checks are deterministic.
Receiver credit bounds accepted new data, while retransmission and probe rules
prevent feedback loss from causing permanent protocol deadlock. Exactly-once
application delivery is limited to a live session; durable crash semantics are
not provided initially.

All feedback and custom metadata travel through the simulated network. Headers,
ACKs, probes, and learned symbols consume bandwidth and face delay and loss.
Start with typed metadata before introducing bounded discrete learned messages.
Metadata may describe local arrivals or buffer state, but cannot reveal hidden
router state. RTT sampling uses echoed timestamps without assuming synchronized
clocks. The first representation is an engineering scaffold and may change as
the learned protocol expands.

\subsection{Observations and architecture}
Initial observations aggregate 20\,ms windows: sent bytes, newly acknowledged
unique bytes, retransmissions, estimated losses, RTT statistics and trends,
outstanding data, pacing rate, congestion window, receiver credit, unsent bytes,
feedback age, previous actions, role, and received metadata. Measurements have
validity masks and fixed-scale normalization. No true capacity or queue occupancy
is supplied to the actor.

The starting model is a causal transformer with $H=64$ windows, two layers,
width 64, four attention heads, and actor and value heads. These are proposed
hyperparameters, not validated selections. Model comparisons include a stacked
history MLP and a recurrent policy at comparable budgets. Session resets clear
history, while rollout boundaries retain the required preceding context.

The implementation also supports a flattened-history MLP and a GRU over the same
masked finite window. A 16-window transformer of width 64 with two layers has
102,723 parameters; an MLP of width 152 with three layers has 103,211; and a
two-layer GRU of width 104 has 105,875. These match parameter budgets rather than
computation. The GRU carries no additional hidden context between policy calls.
Inactive training endpoint slots are excluded from the optimization objective.

\subsection{Action spaces}
The first policy controls pacing with a scalar $a_r\in[-1,1]$:
\begin{equation}
 r_{t+1}=\operatorname{clip}\bigl(r_t\exp(0.2a_r),r_{\min},r_{\max}\bigr).
\end{equation}
A second stage introduces $a_w\in[-1,1]$ with an analogous update for the
congestion window. Bounds do not depend on hidden capacity. Pacing applies to
retransmissions as well as new data; receiver credit constrains newly introduced
payload separately.

Later policy heads choose eligible retransmissions, packet sizes, feedback
timing, bounded metadata, scheduling, redundancy, and available paths. Masked
actions invoke defined protocol operations. They cannot alter payload bytes or
acknowledge data that has not arrived. We introduce heads separately so their
benefit and overhead can be measured. The full system jointly trains sender and
receiver policies with explicit roles.

\section{Reference environment and GPU scaling}
The CPU reference simulates a finite FIFO bottleneck with byte-based
serialization, forward and feedback paths, propagation delay, jitter, and
independent random and queue-overflow loss. Scripted cases add duplication,
reordering, burst loss, changing capacity, and temporary disconnections. A
20\,ms control interval does not quantize all network events to 20\,ms: the
reference preserves event timing and paced departures.

Proposed initial ranges are 0.5--20\,Mbit/s, 10--200\,ms base RTT,
0--5\% random loss, and 8--128\,KiB bottleneck queues. Begin with one transfer,
then add 2--8 competing senders and receiver-side control. Short text checks
correctness; longer or repeated messages exercise steady-state congestion.
Large messages use bounded buffering and backpressure.

The GPU backend will batch independent networks with interacting endpoints
within each network. Packet slots, queues, histories, and actions have fixed
bounded layouts and explicit masks. Immutable payload references avoid repeated
string copies. Physical queue drops and implementation storage exhaustion are
separate outcomes. A bounded event or fine-tick implementation must be checked
against CPU timing and conservation invariants before training at scale.

Benchmark batches of 128, 512, and 2,048 four-sender networks without presuming
they fit a particular device. Keep simulation, inference, rollouts, and updates
on-device. Store observations once and reconstruct history windows for training.
Time remains causal; parallelism comes primarily from independent networks.
Report memory, packet events per second, transitions per second, simulated time
per wall-clock second, and learning time to a stated target. Multi-GPU sharding
is deferred until a single-device implementation is validated.

\section{Training objective and evaluation plan}
The initial PPO objective rewards unique acknowledged payload progress $B_t$
and penalizes estimated RTT inflation $D_t$ and overhead wire bytes $O_t$:
\begin{equation}
 R_t = \frac{B_t}{B_0}-\alpha\frac{D_t}{D_0}
       -\beta\frac{O_t}{B_0}.
\end{equation}
Here $B_0,D_0$ are fixed reference scales. Overhead includes feedback,
headers, probes, metadata, and duplicate payload. Finite transfers also receive
a verified completion bonus and an unfinished-transfer penalty. When feedback
becomes learned, training credit comes from independently verified receiver
delivery, rather than trusting generated acknowledgments; that verifier does
not become an actor observation. Delayed feedback must be considered in credit
assignment.

A second, versioned completion-focused objective addresses the original short
discount horizon. For message length $L$, active elapsed time $\Delta t$, RTT
inflation $D_t$ in seconds, progress $B_t$, and overhead $O_t$, define
\begin{align}
 R_t^{(2)} ={}& \frac{B_t}{\max(L,1)}
 + 2\,\mathbf{1}_{\mathrm{complete}} - 2\,\mathbf{1}_{\mathrm{timeout}}
 - 0.05\Delta t \nonumber\\
 &-0.01\frac{D_t}{0.1}\Delta t
 -0.1\frac{O_t}{\max(L,512)}.
\end{align}
The original profile uses $\gamma=0.99$ and the second $\gamma=0.999$. At
20\,ms intervals, a terminal reward after 20 seconds is weighted by approximately
$0.000043$ or $0.368$, respectively. The second profile gives waiting an explicit
cost and retains more of the completion signal. Evaluation specifies one common
profile for all controllers; different reward definitions are not aggregated.
Neither profile relaxes byte correctness.

The robust training sampler includes one to four senders, up to 15\% random loss
and 2\% corruption, changing capacity, asymmetric feedback capacity, and outages.
Clean episodes remain in the mixture. This sampler is separate from the frozen
evaluation manifest. The experiment compares two training seeds per architecture
and an old-reward ablation, selects on validation seeds, then evaluates fresh
final-test seeds.

For joint multi-flow training, evaluate a network objective such as
$\sum_{i\in\mathcal{A}_t}\log(\epsilon+g_t^i/g_0)$, where
$\mathcal{A}_t$ contains active backlogged flows and $g_t^i$ is verified goodput.
Declare how joins, leaves, and idle periods affect the active set. Network-level
training signals are distinct from runtime peer communication. A centralized
critic, if used, must be reported explicitly.

\begin{table}[ht]
\centering
\small
\begin{tabular}{@{}p{0.23\linewidth}p{0.69\linewidth}@{}}
\toprule
Question & Planned evidence \\
\midrule
Correct delivery & Exact bytes, empty and Unicode messages, duplicated packets,
lost feedback, outages, deadlines, and bounded-buffer behavior. \\
Control quality & Completion-time distributions, goodput, RTT percentiles,
wire overhead, and recovery traffic. \\
Fairness & Per-flow outcomes, Jain's index, and worst-flow waiting time under
joins, leaves, and heterogeneous RTTs. \\
History benefit & Transformer versus MLP and recurrent models; history-length
and inference-cost ablations. \\
Protocol freedom & Sender-only versus joint training; metadata budgets,
feedback timing, and recovery-action ablations. \\
Generalization & Held-out parameter combinations, traces, traffic mixes,
and multiple independent training seeds. \\
Scalability & Reference agreement, timing error, memory, environment throughput,
and training time to target. \\
\bottomrule
\end{tabular}
\caption{Evaluation plan beyond the preliminary software diagnostics.}
\end{table}

Fixed-rate and AIMD baselines use identical reliability code and paired scenario
seeds. AIMD alone is not a full TCP implementation. Later comparisons use actual
TCP or QUIC controllers under a common emulator. A single-message UDP adapter
reuses the protocol and policy interface. It saves received text atomically,
verifies a completion digest, and retains receiver state for a finite interval
to repeat lost final feedback. A worker thread computes policy decisions while
the sender continues under the last action, following the asynchronous execution
concern raised by MVFST-RL~\cite{sivakumar2019mvfst}. Multi-host performance remains
to be measured under controlled impairment.

\section{Preliminary software experiment}
An exploratory comparison trains two seeds per history architecture for 400
updates under the robust sampler and completion-focused reward. Each run collects
approximately half a million active transitions. A seventh run trains a robust
transformer with the original reward as an ablation. Validation uses ten paired
seeds per fixed scenario, and selects first by completed transfers and then by
mean per-scenario completion-focused return. Final evaluation uses twenty fresh
seeds per scenario, after recording the selection.

\begin{table}[ht]
\centering
\small
\begin{tabular}{@{}lrrrr@{}}
\toprule
Scenario & Fixed & AIMD & Previous & Selected \\
\midrule
Clean & 40/40 & 40/40 & 40/40 & 40/40 \\
Lossy/corrupt & 40/40 & 40/40 & 40/40 & 40/40 \\
Changing/outage & 40/40 & 40/40 & 38/40 & 40/40 \\
Four-sender stress & 80/80 & 4/80 & 26/80 & 80/80 \\
\bottomrule
\end{tabular}
\caption{Completed exact transfers on reserved final-test seeds. ``Previous'' is
the earlier 300-update sender; ``Selected'' is transformer training seed 43.}
\end{table}

The selected transformer completes all 200 transfers without wrong payloads, as
does fixed-rate. This shows progress over the previous learned sender, not general
superiority over conventional transport. Seed sensitivity is substantial: the
other completion-focused transformer completes none of its 100 validation
transfers, while both MLPs and both GRUs complete all of theirs. The old-reward
ablation also completes validation, so the experiment does not isolate the new
reward as the cause of improvement. Scenarios are fixed diagnostic conditions,
not an Internet deployment benchmark; flows within a trial are dependent.

\section{Receiver-only diagnostic}
We freeze the AIMD sender and train a separate eight-step history transformer
(width 32, one layer) on 22 receiver-local features. Two actions select an ACK
batch target from $\{1,2,4,8\}$ and a maximum delay in $[0,25]$ ms. A timer flushes
partial batches; completion feedback is immediate. An explicit experimental
packet kind carries bounded batches of actual received ranges and transmission
IDs. Encoded feedback consumes reverse-link capacity, and explicit ACK delay is
removed from RTT samples. This extension is simulator-only.

Three training seeds receive 80 episodes each. Validation seeds 1000--1004 select
the learned candidate and a deterministic baseline by completion, then common
undiscounted transport-v2 return. Test seeds 2000--2019 remain separate. Six
scenarios include clean, two asymmetric reverse paths, loss, changing forward
capacity, and an outage; the outage family is absent from training. Two flows
share each trial. Other scenario families overlap training, so held-out seeds do
not constitute general out-of-distribution evaluation.

The diagnostic win requires exact completion of every transfer for both selected
policies, plus either lower elapsed time with a paired 95\% interval below zero
and no greater mean feedback, or lower feedback with a paired interval below
zero and elapsed-time upper confidence bound within $+5\%$. Bootstrap resampling
uses whole seeds, retaining dependent flows and scenarios together.

Validation selects learned seed 43 and deterministic batch-8/25-ms feedback.
On 240 final-test transfers, the learned receiver completes 223, immediate ACKs
complete 215, and the selected deterministic baseline completes 240. No wrong
payloads occur. Feedback totals are respectively 3,099,580, 3,765,162, and
2,505,897 bytes. The learned candidate has 23.7\% more feedback and 15.3\% greater
mean elapsed time than the selected baseline, counting deadlines for failures.
All 17 learned-candidate failures occur with the 24-kbit/s reverse path.
The win condition fails: batching helps this testbed, but this short learning
experiment does not outperform tuned deterministic batching. Simulated elapsed
time excludes actor inference cost. Configurations, raw paired records, training
logs, and a separately versioned candidate checkpoint accompany the repository.

\section{Receipt recovery and persistent degradation}
An opt-in receipt packet repeats the contiguous received prefix and up to eight
selective byte intervals. A digest-verified full prefix confirms the whole
message even when previous feedback was lost. Only explicitly echoed
transmissions yield RTT samples; summaries cannot credit unsent segments.
The first completion receipt is immediate, while duplicate completion responses
follow bounded scheduling. The socket adapter supports receipt summaries, but
learned receiver scheduling remains simulator-only.

We compare preset training with persistent randomized degradation and a
one-observation ablation, each with three seeds and 80 episodes. Randomized
episodes contain a gradual decline, a sustained poor interval, and recovery;
capacity, residual loss, and access delay covary. Queueing remains
traffic-dependent. Named Wi-Fi cases cover stable service, contention, mobile
fading, and clustered loss. These are uncalibrated service models above the MAC
layer, not an implementation of 802.11. Policies never observe hidden schedules.
Randomization varies one to four flows; equal episode budgets therefore do not
imply equal transition or compute budgets.

Validation seeds 3000--3004 select candidates. On diagnostic seeds 4000--4019,
preset/history-8, randomized/history-8, and randomized/history-1 candidates
complete 239/240, 238/240, and 238/240 transfers; deterministic batch-8/25-ms
feedback completes 240/240. Holding the original receiver-v1 weights fixed,
switching only the protocol raises completion from 223/240 to 239/240 on these
same fresh seeds. This isolates a recovery benefit from retraining.

On a separate 560-transfer generalization suite, selected using no outcomes from
seeds 5000--5019, the three learned candidates complete 551, 551, and 552
transfers; the deterministic baseline completes 553. Their Wi-Fi subset results
are 159/160, 160/160, and 160/160, versus 160/160 for the deterministic baseline.
No wrong payloads occur. The suite includes longer RTTs, more flows, larger
messages, smaller queues, stronger bursts, and reverse capacity beyond training
support. All learned candidates fail the diagnostic win condition. Neither
longer neural history nor domain randomization shows an overall advantage at
this budget. The one-observation ablation still receives endpoint counters and
timers; it is not a host without memory. Real-device traces are needed to
calibrate persistence and assess transfer beyond these stylized models.

\section{Control headroom and feedback-wait budgets}
A wider study first sweeps constant ACK controls. Selecting a constant separately
for each scenario using validation labels reduces test elapsed time by 4.9\%
but increases feedback by 25.6\%; this is neither deployable nor an upper bound
on a time-varying policy. A categorical receiver initialized at the strong
baseline, trained for 320 episodes per seed, selects its initial checkpoint on
validation and exactly matches the deterministic control on 240 diagnostic
transfers. Increasing training alone does not establish an advantage.

The old action range also caps ACK waiting at 25 ms, while sender timers do not
explicitly budget intentional feedback waiting. We configure a 250-ms maximum
feedback-wait budget at both endpoints and add that budget to RTO. This is a
shared simulator configuration, not on-wire negotiation. Validation seeds
17000--17004 select batch-8/250-ms feedback; test seeds 18000--18019 remain
separate. All three diagnostic arms complete 240/240 exact transfers:
\begin{center}
\begin{tabular}{lrr}
\toprule
Protocol arm & Mean seconds & Total wire bytes \\
\midrule
Old ACK policy and timers & 2.367 & 15,608,314 \\
Old ACK policy, budgeted timers & 2.172 & 12,328,675 \\
Longer ACK wait, budgeted timers & 1.491 & 10,558,204 \\
\bottomrule
\end{tabular}
\end{center}
The combined change reduces mean elapsed time 37.0\%, ACK bytes 70.5\%, and
total wire bytes 32.4\%. The paired seed-block 95\% interval for elapsed-time
difference is $[-1.067,-0.675]$ seconds. On generalization seeds 19000--19019,
completion rises from 551/560 to 560/560. The timer-only arm is faster there,
while the longer-wait arm saves more feedback: these are distinct tradeoffs,
and generalization elapsed means include deadlines for failures.

These gains are deterministic protocol improvements. A categorical policy over
21 feedback controls, trained with the new budget, completes 240/240 diagnostic
transfers but is slower and uses more ACK bytes than its validation-selected
constant baseline. On another generalization split it completes 559/560,
compared with 560/560 for the latency-focused deterministic control. A learned
sender correction likewise has no clear total-wire advantage over strong
deterministic senders. A feedback-only criterion would conceal extra data
traffic; the sender analysis was corrected after inspecting its breakdown and
is exploratory. Finally, a new mixed workload combines two 128-KiB bulk
transfers with sixteen staggered 1-KiB messages. Frozen controllers expose bulk
timeouts and multi-second short-message tail latency, motivating future
application-aware training rather than inferring superiority from small,
simultaneous-transfer diagnostics.

\section{Persistent-session joint experiment}
The joint implementation uses two local transformer actors and a training-only
centralized critic. Sender actions control rate, flight window, and queued-message
selection; receiver actions control bounded ACK batching and timing. A fixed
reliability kernel enforces exact binary recovery. Histories persist across
messages. Interactive text has a 100-ms deadline; images and map deltas provide
bulk progress objectives. Backpressure counts blocked interactive offers as misses.

The version-two curriculum samples correlated fades, contention, outages,
asymmetric service, and feasible through overloaded wire loads. It accounts for
headers and bidirectional feedback; an earlier undercounting sampler was retired.
Shared airtime approximates Wi-Fi service above the MAC and is not a calibrated
802.11 model. Six five-million-decision runs cover three joint seeds, one-step
history, and frozen-role ablations. CPU event workers feed batched GPU inference
and optimization. Actor exports exclude the centralized critic.

Validation chooses among deterministic rate, scheduling, and feedback controls
before disjoint held-out testing. Acceptance requires a 20\% relative deadline-miss
reduction with a paired bootstrap interval below zero, at least 95\% of reference
image/map goodput and delivered volume, no completion-ratio loss, at most 105\%
wire cost per delivered byte, and zero corrupt delivery. Training reward alone
cannot establish this result. The binary UDP adapter optionally runs a local
learned sender with deterministic remote feedback; deploying a learned receiver
and persistent joint sessions on hardware remains future work.

\subsection{Selection and held-out result}
Calibration uses seeds 30000--30003, checkpoint screening 40000--40003, and full
validation 41000--41011. Joint training seeds are 42, 43, and 44; one-step history,
sender-only training, and receiver-only training use seed 42 and the same nominal
five-million active-endpoint decision budget. Full validation compares three initial snapshots with six trained candidates
chosen by screening.
All initial snapshots implement the same deterministic behavior. Validation
selects \texttt{seed42/initial.pt}; training does not improve the selected policy.
The deterministic reference is fixed-rate/EDF with batch-8 feedback and a 250-ms
maximum wait. The initial actors follow fixed-rate/EDF and the 25-ms feedback
teacher. Their difference is a deterministic feedback configuration, not learning.

Fresh test seeds 50000--50039 and OOD seeds 60000--60039 are disjoint from
selection data. Each case has 20 seconds of production and up to 20 seconds of
draining. Table~\ref{tab:joint-primary} reports the completed primary comparisons.
Goodput counts receiver-completed object bytes divided by total simulated session
time; neural execution time is excluded from simulation. OOD includes asymmetric
handoffs, longer RTTs, and additional hosts, including physically infeasible
100-ms intervals. Desired interactive offers, including blocked offers, remain
in the deadline denominator.

\begin{table}[ht]
\centering\small
\begin{tabular}{llrrrr}
\toprule
Suite & Control & Misses (\%) & Image MiB/s & Map MiB/s & Wire/useful \\
\midrule
Test & Reference & 40.71 & 0.0758 & 0.0760 & 1.569 \\
Test & Selected initial & 40.00 & 0.0795 & 0.0773 & 1.613 \\
OOD & Reference & 79.68 & 0.0421 & 0.0320 & 1.839 \\
OOD & Selected initial & 79.07 & 0.0486 & 0.0381 & 1.823 \\
\bottomrule
\end{tabular}
\caption{The selected joint architecture uses its untrained baseline-following
actors. This table does not compare a successful learned policy against a baseline.}
\label{tab:joint-primary}
\end{table}

Test relative miss reduction is 1.74\%, with a paired seed-block bootstrap
95\% interval for the absolute miss-rate difference of
$[-1.729,0.194]$ percentage points. The interval includes zero and the reduction
is below 20\%. The strict confirmed-completion-ratio guard also fails:
99.96974\% versus 99.97109\%. Image/map goodput and volume, wire-cost, and
integrity guards pass. OOD relative reduction is 0.76\%, with a difference
interval of $[-0.947,-0.348]$ percentage points; this is below the target despite
excluding zero. No learning win is established in either suite. The reliability
kernel reports no wrong delivery in the completed cases. Cross-play and
separately trained role/history ablations are reported alongside the primary
policy, not used to change the frozen test selection. All 560 case/mode evaluations
finish (seven modes, 40 cases, two suites), with no wrong delivery observed.

\paragraph{Runtime and reproducibility limits.}
Training continued from saved actor, critic, optimizer, RNG, and decision-counter
states after interruptions. Collector worlds and histories restarted with fresh
episode seeds, so this is not bit-for-bit uninterrupted training. The history-one
run returned native exit code 139. Later worker diagnostics captured a segmentation
fault in CPython 3.12.10's JSON encoder during receipt serialization at network
seed 50016. A Python-level container encoder preserves canonical wire bytes in
parity tests and avoids that observed path; the underlying native cause remains
unconfirmed. Per-case persistence protects completed held-out sessions.
Asynchronous independent-world evaluation avoids global per-tick worker barriers;
GPU batching may change floating-point round-off and is not claimed bit-identical
across execution configurations. Runtime notes, source/configuration hashes,
checkpoints, and negative results accompany the experiment register in
\texttt{docs/experiments/register.md}.

\paragraph{Next throughput objective (planned).}
A separate version-two experiment will keep bulk traffic continuously backlogged
and maximize verified useful bytes per second, including finite-object completion
time. Mixed traffic retains the 100-ms interactive target through explicit
noninferiority constraints on deadline misses, bounded queues, wire overhead, and
active-host service. Retransmissions and dropped bytes are not useful throughput.
The proposed acceptance threshold is 10\% more goodput with a positive paired
confidence interval, subject to latency and reliability guards and comparison
against a strong adaptive deterministic reference. This objective requires new
training and evaluation; it does not redefine the current deadline-first study.

\paragraph{Bounded device checks.}
A desktop sender transferred synthetic Unicode text, valid 256-by-256 RGB PNG
images, and 512-KiB map buffers to Herman over stationary Wi-Fi and Caliban over
Ethernet. Three trials per payload and device succeeded with deterministic
control (18/18) and with the separately trained sender-only actor and a
deterministic receiver (18/18). The latter adapter constructs local sender
histories and controls rate/window; a single pending message has trivial
scheduling. This differs from the persistent mixed training workload and does
not deploy a learned receiver. The sequential, nonrandomized checks establish
byte-exact recovery, not a learned performance advantage. Eighty ICMP probes per
recording provide RTT/loss observations, not available-bandwidth estimates.

\section{Implementation stages and open questions}
Implemented stages include reliable CPU transfer, history-based sender learning,
receiver learning, joint sender/receiver training over CPU packet workers, and
bounded binary datagram checks on two real devices. A GPU-native packet backend,
joint learned-receiver hardware deployment, and maximum-goodput training remain
future work.
The full target combines these pieces with reproducible configurations,
checkpoints, bounded resources, and explicit failure semantics.

Open questions include the value of attention versus cheaper memory models,
feedback representations that justify their wire cost, generalization under
unseen workloads, and the fidelity needed for simulator-trained policies to
transfer to real packet timing. Arbitrary metadata and self-play enlarge the
design space but do not themselves establish better communication behavior.

\bibliographystyle{plain}
\bibliography{references}
\end{document}
